\chapter{Literature Review}
\label{chap:lr}

\section{Questions}

Based on my goal, I have formulated the following questions for the literature review:

LRQ\#1: What does the scientific community understand by branch-and-bound and backtracking methods?
What are the parts common to all branch-and-bound-based and backtracking-based algorithms?
Which parts of the the algorithms vary from problem to problem? \\

LRQ\#2: What were the previous attempts to formalize branch-and-bound and backtracking methods
and algorithms based on them? \\

LRQ\#3: What are ways in which these algorithms can be or are commonly implemented? \\

\section{Searching procedure}

For each question, I searched for articles in ScienceDirect,
Google Scholar, and ResearchGate, combining the following keywords to form search queries
for my questions:

\begin{itemize}
    \item Backtracking (method)
    \item Branch and bound (method)
    \item General / Generic (LRQ\#2)
    \item Formal / Formalization / Verification (LRQ\#2)
\end{itemize}

\section{Answers to the questions}

\subsection{LRQ\#1: Understanding backtracking and branch-and-bound}

\subsubsection{Branch \& Bound method}

In the most general case studied in \cite{clausen1999principles}, \textbf{Branch \& Bound (B\&B)}
algorithms solve optimization problems, that is, minimize an \emph{objective function} in the
\emph{feasible solution space}, which is limited by constraints. To do this, B\&B algorithms
implement a search of a "dynamically developed" (lazily-generated) graph
(also referred to as \emph{virtual graph}
    \cite{shilov2011rutemplateverif} \cite{shilov2012verifmono}),
which acts as a search tree, where every node represents a subproblem derived from the
original problem through addition of constraints.

Initially, the tree consists of just the root node, which represents the original problem.
On every iteration, a node is extracted from the set of \emph{live nodes}.
It is then evaluated: the subproblem is broken down to
two \cite{narkawicz2013formalnasa} or more \cite{clausen1999principles}
subproblems according to a
\emph{branching rule}, new nodes become children of the evaluated node and are added
to the set of live nodes. The \emph{bounding function} is evaluated for nodes in
consideration: if the value of the bounding function at the candidate node is greater than
or equal to (i.e., not better than) the value of the objective function at the current
best solution (\emph{incumbent}), the candidate can be discarded without descending into
its subtree, as its solution space can not contain a solution better than the incumbent
(due to properties of the bounding function, discussed below).

When a subproblem can not be broken down further, the objective function is evaluated at
it and compared to the objective function at the incumbent, and the best one is kept.

\subsubsection{Problem-specific parts of B\&B}

While the core part of walking the search tree is common to all the B\&B algorithms,
there are parts that depend on a concrete problem that an algorithm is solving.

For clarity, I will discuss these parts along with the example of the Sudoku solving
B\&B-based algorithm \cite{indriyono2024sudoku}. Every node of the solution tree
in the example represents a partially filled Sudoku board.

Following the terminology of \cite{clausen1999principles}:

\begin{itemize}
    \item \emph{Objective function} is a function of a feasible solution. The algorithm shall
        minimize it.
        Objective function can be evaluated for a concrete solution
        (represented in the tree with a leaf node) rather than a solution space. \\
        As in Sudoku there is no "worse" or "better" solutions, only valid or invalid ones,
        we could define a fairly straightforward objective function that is equal to
        $1$ for a valid solution and $2$ for an invalid solution (as less is better).

    \item \emph{Root node} represents the original problem to be solved. \\
        In the Sudoku example,
        the root node represents the pre-filled Sudoku board, as given initially;
        pre-filled values are the initial constraints.

    \item \emph{Branching rule} is a rule of (function that defines) how to break down a subproblem
        into further subproblems (i.e., generate descendants of a node). \\
        The following shall hold for a branching rule to be valid:
        \begin{enumerate}[(a)]
        \item solution space of the parent node shall be equal to the union of solution spaces of its
            direct descendants, and
        \item solution spaces of sibling nodes shall not intersect (i.e.,
            constraints imposed on different descendants shall be mutually exclusive).
        \end{enumerate}

        In the Sudoku example, one branching rule idea is to select an empty cell
        and try to insert every possible digit (1-9) in it, thus creating 9 descendants of a node.
        If the Sudoku board in a node does not have any empty cells, no further branching is possible,
        and that node is a leaf node.

    \item \emph{Bounding function} is a computationally-cheap function of a solution subspace
        that estimates the best potentially achievable solution in that
        subspace. \\
        The following shall hold for a bounding function $g$ to be valid:
        \begin{enumerate}[(a)]
            \item Bounding function is no worse than the objective function $f$:
                \[
                g(\{X\}) \leq f(X) \:\text{for any feasible solution}\: X
                \]
            \item Narrowing solution space down shall not improve the best solution estimate:
                \[
                g(S) \leq g(S') \:\text{for}\: S \supseteq S'
                \]
                or, equivalently
                \[
                \forall S, T: g(S \cup T) \leq \min(g(S), g(T))
                \]
        \end{enumerate}

        In the Sudoku example, the bounding function will be used merely to discard subtrees
        in which we can never find a solution (i.e., rule out partial solutions where the
        pre-filled part is already invalid). Thus, we could define
        the bounding function to be equal to $1$ if the partially filled board is valid and $2$
        otherwise.

        In the Sudoku example, the bounding function happens to be very similar to the objective
        function; in more advanced optimization problems, a more sophisticated bounding
        function would be useful.

    \item \emph{Selection strategy} is a strategy according to which a live node is selected
        at the beginning of every iteration (a-la the walk order of the virtual graph). \\
        Common selection strategies are:
        \begin{itemize}
            \item Breadth-first search (BFS) evaluates nodes of the tree layer by layer;
            \item Best-first search (BeFS) greedily selects a node with the lowest value of the
                bounding function out of the live nodes;
            \item Depth-first search (DFS) walks the tree in a DFS manner.
        \end{itemize}

        For many optimization problems, it makes sense to use BFS or BeFS to arrive at the best
        solution earlier (discard more solution spaces without descending into subtrees).
        For the Sudoku example, however, it would not benefit the performance, as the bounding
        function conveys too little information about the solution subspace, so it is more efficient
        to use DFS as the most memory-efficient option.

    \item \emph{Evaluation strategy}, either \emph{eager} or \emph{lazy}, affects when the
        algorithm compares candidate nodes against the incumbent to potentially discard them:
        \begin{itemize}
            \item With \emph{eager} evaluation, a selected live node is first \emph{branched}
                (i.e., split into subproblems), then
                new nodes are compared against the incumbent and are only inserted into the set
                of live nodes if they have the potential to produce a better solution
                (discarded otherwise);
            \item With \emph{lazy} evaluation, a node selected from live nodes is first compared
                with the incumbent (and can be discarded at this point), then branched,
                and the descendant nodes are inserted to the set of live nodes unconditionally.
                \todo{(Nechaev) but what if we combine eager+lazy into an "aggressive" strategy:
                    evaluate both on insert and on extraction? - (Shilov) It makes sense, of course, but should be named "joint" rather than "aggressive".}
        \end{itemize}
\end{itemize}

\subsubsection{Backtracking and other special cases}

% TODO: LIE LIE LIE LIE
% TODO: clausen1999principles does not mention backtracking at all!!!!!
According to \cite{clausen1999principles}, \textbf{Backtracking} is a special case of the
Branch \& Bound method with DFS used as the selection strategy.

Brute-force search for optimization problems could be seen as another special case of the
Branch \& Bound method with the bounding function $g(x) = -\infty$.
Indeed, in this case, all nodes are evaluated
and no subtree is ever discarded, leading to the exhaustive search through all the possible solutions
(leaf nodes).

\subsection{LRQ\#2: What were the previous attempts to formalize branch-and-bound and backtracking methods and algorithms based on them?}

In \cite{narkawicz2013formalnasa}, authors present and formally verify a "generic branching
algorithm\footnote{The terms \emph{generic algorithm} and \emph{template} are used interchangeably
                    here and after}
for solving global optimization problems". Their algorithm template is to be instantiated with
functions (such as the objective function, the bounding function, etc) to form a concrete
algorithm for solving a concrete optimization problem, which would be easier to formally verify
because the core (generic) part of it is proven correct as long as supplied functions
satisfy certain criteria.

However, because the template in \cite{narkawicz2013formalnasa} was designed specifically for
numerical optimization problems, it does not fit the most general B\&B use case for the following
reasons:

\begin{enumerate}
    \item Because the researchers' goal was to simplify the verification of algorithms instantiated
        from their template, the template is specialized for numerical optimization. It moves the
        burden of verifying common numerical-related statements to the authors of the generic
        part of the algorithm, however, it makes it harder to use their template for other kinds of
        optimization problems. For example, the template in \cite{narkawicz2013formalnasa}
        is designed to branch in a way that only makes sense when the objective function is seen as
        a numerical function of multiple variables, each having a 1-dimensional domain.
    \item Some aspects of the branch-and-bound technique, such as evaluation strategy
        and selection strategy, can not be specialized in the
        generic algorithm from \cite{narkawicz2013formalnasa} because they are hard-coded into it.
\end{enumerate}

The work in \cite{narkawicz2013formalnasa} further continues in \cite{smith2015rigorous},
and is implemented as a C++ library called
Kodiak \footnote{Electronically available at \url{https://github.com/nasa/Kodiak}},
which was not itself formally verified, but, according to the authors, "closely follows
the formally verified depth-first branch and bound algorithm with generic types for problem domains
and solution types" (referring to \cite{narkawicz2013formalnasa}). Being an implementation of the
generic algorithm discussed above, it suffers from the same issues: it is specialized to solve
numerical optimization problems and does not allow to specialize some aspects, such as the
virtual graph walk order.

\subsection{LRQ\#3: What are ways in which these algorithms can be or are commonly implemented?}

There were implementations of both backtracking-based and branch-and-bound-based algorithms
for specific problems
(e.g., \cite{bard1990bilevel}, \cite{breuel2003geometric}, \cite{lalami2012gpu}),
as well as implementations generalized to varying extent
(e.g., \cite{narkawicz2013formalnasa}, \cite{smirnov2017concur}, \cite{finkel1987distrib},
\cite{prenner1972proglangs}, \cite{johnson1988modular}).

In most cases, backtracking is implemented with recursive function calls,
e.g., in \cite{narkawicz2013formalnasa}, \cite{bard1990bilevel},
and branch-and-bound is implemented with a queue or a priority queue,
depending on the selection strategy, e.g., in \cite{breuel2003geometric}.

In special cases, other technical implementations are used, although the computation logic
is roughly the same. For example, in \cite{lalami2012gpu}, authors implement a
breadth-first search with items "sorted according to decreasing profit per weight ratio"
(for most natural definitions of boundary functions, it is the same as best-first-search),
however, because in this work computations are parallelized with a GPU
(with some parts performed on a CPU), the technical implementation is somewhat more
sophisticated than a mere loop on a priority queue.

Distributed branch-and-bound search is another example in which the technical implementation
becomes more complicated, although it follows the same general model. For example,
in \cite{smirnov2017concur}, worker machines are assigned subtrees of the main search tree;
during the computation, they only exchange their incumbent values. In particular, while solving one
problem, workers can use different selection strategies, which can improve overall performance.
In \cite{finkel1987distrib}, a different distributed algorithm is presented: the branching and
node processing stages are separated, and worker machines exchange information about their
subtrees and can be dynamically assigned nodes to process.

There were also examples of more abstract research. For example, the
authors of \cite{prenner1972proglangs} studied the possibilities of implementing
backtracking in various programming languages, trying to stay as decoupled from specific language
features as possible. This work, however, seems to focus on programming language design
rather than backtracking search algorithms. In addition to that, the paper is relatively old and
only presents theoretical ideas of backtracking implementation. Due to these reasons, it is unlikely
that the article will be useful in the presented work.

Finally, \cite{johnson1988modular} presents a modular implementation of backtracking and
branch-and-bound: core functions common to many concrete algorithms are implemented separately
from problem-specific parts. The authors analyze the ease of implementation, as well as performance
of assembly line balancing algorithms based on modular implementation compared to direct
implementation. The authors also devise two different kinds of DFS branch-and-bound algorithms,
referred to in the paper as \emph{deep-sea-troll} search and \emph{laser} search. While both
being LIFO search strategies, selecting the right one of the two for a concrete  problem may
affect performance (see \cite{johnson1988modular} for details).
% It is, however, possible to encapsulate the choice
% between deep-sea-troll and laser search strategies in the branching function, as long as the
% collection of nodes that it returns is ordered: if the core library implements the laser search,
% then a branching function returning new nodes in arbitrary order makes it a true laser search,
% while a branching function that returns its values in order of decreasing profit will turn shall
% result in an algorithm analogical to a deep-sea-troll search.

\section{Conclusion}

The relevant literature provides definitions and descriptions of the backtracking and
branch-and-bound methods (both specialized for particular problems and general ones) and
examples of their application to optimization problems.

I have found theoretical research studying implementation
of branch-and-bound for concrete problems
\cite{clausen1999principles}, \cite{bard1990bilevel}, \cite{indriyono2024sudoku},
as well as works that implement and/or verify versions of the branch-and-bound algorithm
generalized to some extent
\cite{narkawicz2013formalnasa}, \cite{johnson1988modular}
or designed for special computation environment, e.g., GP$^2$U \cite{lalami2012gpu} or
distributed versions of the algorithm \cite{finkel1987distrib}, \cite{smirnov2017concur}.

However, only in \cite{johnson1988modular} did the authors present an approach that is fully
general both in the problem to solve (as opposed to solutions designed for solving, e.g., global
optimization problems, such as \cite{narkawicz2013formalnasa}) and in the computation method
(as opposed to solutions designed for, e.g., distributed computation, such as
\cite{finkel1987distrib}, \cite{smirnov2017concur}).

In the presented work, I
implement a template library in the Rust programming language that does not rely
on a special computation environment and is general enough to allow for solving
arbitrary problems where branch-and-bound and/or backtracking algorithms can be applied.

% =========== TEMPLATE STUFF BELOW ===========

%\chaptermark{Second Chapter Heading}


% \Blindtext[2]
%
% \section{Another Section}
% \Blindtext[1]
% \newpage
% With the widespread of computing systems, information processing, and net working, the practice of replacing paper documentation to electronic documentation has become more and more common. Electronic documentation within the workplace has several advantages over the traditional one, such as being easy to share, copy and edit the document. However, these advantages also present a problem. A malefactor, having access to the system, can easily copy and leak the document, without leaving any trace. Such actions are virtually undetectable in most systems, so, the malefactor goes unpunished. In this thesis, we propose one solution to the problem: digital watermarking.
%
% Every electronic document within the protected system is marked with an invisible digital watermark, containing information about the user, accessing this particular document. Therefore, in case the protected company discovers the leaked document, they will be able to identify the machine of the malicious person and time when the document was leaked. This will allow inflicting punishment on the malefactor, recovering the costs of the leak, and potentially preventing future ones.
%
% \begin{longtable}{c|c}
% \caption[This is the title I want to appear in the List of Tables]{Simulation Parameters} \label{table:secsimulation_params} \\
% \hline
% A & B  \\
% \hline
% \endfirsthead
% \multicolumn{2}{@{}l}{} \\
% \hline
% A & B \\
% \hline
% \endhead
% \hline
%  \textbf{Parameter} & \textbf{Value}\\
%  \hline
%  Number of vehicles & $|\mathcal{V}|$\\
%  \hline
%  Number of RSUs & $|\mathcal{U}|$\\
%  \hline
%  RSU coverage radius & 150 m\\
%  \hline
%  V2V communication radius & 30 m\\
%  \hline
%  Smart vehicle antenna height & 1.5 m\\
%  \hline
%  RSU antenna height & 25 m\\
%  \hline
%  Smart vehicle maximum speed & $v_{max}$ m/s\\
%  \hline
%  Smart vehicle minimum speed & $v_{min}$ m/s\\
%  \hline
%  Common smart vehicle cache capacities & $[50, 100, 150, 200, 250]$ mb\\
%  \hline
%  Common RSU cache capacities & $[5000,1000,1500,2000,2500]$ mb\\
%  \hline
%  Common backhaul rates & $[75, 100, 150]$ mb/s\\
%  \hline
% \end{longtable}
%
% \begin{figure}[hbt]
% \centering
% \includegraphics[]{figs/IU.png}
% \caption{One kernel at $x_s$ (\emph{dotted kernel}) or two kernels at
% $x_i$ and $x_j$ (\textit{left and right}) lead to the same summed estimate
% at $x_s$. This shows a figure consisting of different types of
% lines. Elements of the figure described in the caption should be set in
% italics, in parentheses, as shown in this sample caption.}
% \label{fig:secex}
% \end{figure}
%
% This description implies several essential properties of the task at hand:
% \begin{enumerate}
%     \item Watermark must contain all necessary information, but still, be placeable and recognizable even on smaller images. The produced watermark must be compact but have the possibility to store enough information.
%     \item To prevent easy tampering, the watermark must be invisible to the naked eye (and, preferably, to basic image parsing tools). If malefactor does not know about the existence of watermark, they might not even try to remove it and disable it.
% \end{enumerate}
